\looseness=-1 Temporal conditioning on step size $\Delta t$ enables consistent gains in model-predictive control, but fails catastrophically in imagination-based actor-critic training: every variable-rate DreamerV3 configuration suffers substantial performance degradation. Factorial diagnostics on two continuous-control tasks identify the combination of trained temporal embedding and multi-rate data exposure as the strongest predictor of this failure, \emph{conditioning-induced value overestimation}, or CIVO, with AUROC${=}1.000$ over $n{=}36$ configurations; cross-environment analysis across seven additional environments confirms specificity to moderate-complexity locomotion. The recurrent state-space model develops a reward prediction gap at $\Delta t{>}1$ that the critic internalizes as optimistic values; the contrast with model-predictive control identifies feedback architecture as a hypothesized structural factor. CIVO exhibits systematic resilience to standard RL remedies: gradient, output, representation, and architecture-level interventions each fail through distinct mechanisms: pathway redundancy, anchor co-drift, phase transitions, and critic convergence in shared-data regimes. Only wide dual-critic combining min-aggregation with increased capacity achieves partial prevention, and open challenges remain in the architectural design of variable-rate model-based RL.
